% Pandoc template for the fed_fraud documents.
% Based on proposal_gen/templates/detailed_proposal_template/detailed_1_classic.tex
% (Classic report): centered title page, table of contents, ruled section
% headings and the Geodesic house footer. Fonts are Latin Modern (lmodern).
% Two layouts: the full report (title page and contents) and, with
% -V brief, a compact one-page-brief header in the same style.
\documentclass[letterpaper,11pt]{article}
\usepackage[T1]{fontenc}
\usepackage[utf8]{inputenc}
\usepackage{lmodern}
\usepackage[letterpaper,top=0.9in,bottom=1.2in,left=1in,right=1in,footskip=0.6in]{geometry}
\usepackage[table]{xcolor}
\usepackage{geodesicbrand}
\usepackage{amsmath,amssymb}
\usepackage{enumitem}
\usepackage{array}
\usepackage{calc}
\usepackage{longtable}
\usepackage{makecell}
\usepackage[most]{tcolorbox}
\usepackage{titlesec}
\usepackage{needspace}
\usepackage{etoolbox}
\usepackage{hyperref}
\hypersetup{colorlinks=true,urlcolor=coolblack,linkcolor=coolblack,filecolor=geoblue,
  pdftitle={Commerce},pdfauthor={\geofirm}}
\urlstyle{same}

\setlist{itemsep=2pt,leftmargin=1.4em}
\setlength{\parindent}{0pt}
\setlength{\parskip}{6pt}
\providecommand{\tightlist}{\setlength{\itemsep}{0pt}\setlength{\parskip}{0pt}}
\setcounter{secnumdepth}{0}

% Section style: house rule under each heading
\titleformat{\section}{\Large\bfseries\color{coolblack}}{\thesection}{0.6em}{}[{\color{coolblack}\titlerule[0.5pt]}]
\titleformat{\subsection}{\large\bfseries\color{geomid}}{\thesubsection}{0.6em}{}
\titleformat{\subsubsection}{\normalsize\bfseries\color{geomid}}{\thesubsubsection}{0.6em}{}
\titlespacing*{\section}{0pt}{10pt}{6pt}

% Tables are written by latex/filters.lua (navy header, banded rows)
\setlength{\LTpre}{6pt}\setlength{\LTpost}{6pt}

\begin{document}

% BRIEF HEADER
\thispagestyle{geodesic}
\begin{minipage}[c]{0.13\textwidth}\includegraphics[width=\linewidth]{icon.png}\end{minipage}\hfill
\begin{minipage}[c]{0.84\textwidth}
  {\textsc{\large \geofirm}}\par\vspace{2pt}
  {\small\color{geogray}Federal fraud expansion\enspace\textperiodcentered\enspace When: \compiletime}
\end{minipage}\par\vspace{6pt}
{\color{coolblack}\rule{\linewidth}{0.3mm}}\par\vspace{2pt}
{\LARGE\bfseries\color{coolblack} Commerce}\par
\vspace{2pt}{\large\color{geomid} One-page brief}\par
\vspace{-2pt}{\color{coolblack}\rule{\linewidth}{0.3mm}}\par\vspace{4pt}
\pagestyle{geodesic}

\textbf{Priority:} 11 of 33 (score 3.3). \textbf{FY2026 outlays:} \$20B
(USAspending). \textbf{As of:} 2026-10-27. Public data only.

\hypertarget{bottom-line}{%
\section{Bottom line}\label{bottom-line}}

Commerce is a small spender today with two very large programs now
paying out: NTIA's \$42.45B Broadband Equity, Access and Deployment
(BEAD) program, passed through states to internet providers, and the
CHIPS manufacturing incentives. Big new money to new or re-formed
recipients is the classic setup for fraud, and both programs publish
recipient-level data. This is a chance to screen before the money is
spent, not after.

\hypertarget{exposure}{%
\section{Exposure}\label{exposure}}

\begin{itemize}
\tightlist
\item
  \$20B in outlays but \$52B in budgetary resources in FY2026,
  reflecting BEAD and CHIPS money not yet spent.
\item
  Broadband subsidy fraud has precedent in the FCC's E-Rate and rural
  programs (see \url{fcc.pdf}): inflated costs, coverage claims that do
  not match service, and bid rigging.
\end{itemize}

\hypertarget{public-data-we-can-use}{%
\section{Public data we can use}\label{public-data-we-can-use}}

\begin{itemize}
\tightlist
\item
  \textbf{State BEAD subgrantee selections} (published by each state)
  and NTIA approvals.
\item
  \textbf{FCC National Broadband Map:} provider-reported coverage by
  location.
\item
  \textbf{USAspending} for CHIPS, EDA and NIST awards.
\item
  \textbf{SAM.gov} and the FCC's USF data for provider histories.
\end{itemize}

\hypertarget{detection-plays}{%
\section{Detection plays}\label{detection-plays}}

\begin{enumerate}
\def\labelenumi{\arabic{enumi}.}
\tightlist
\item
  \textbf{Coverage claims vs.~reality:} BEAD-funded locations already
  shown as served on the FCC map, or providers claiming coverage they
  cannot show.
\item
  \textbf{Provider history:} BEAD winners with prior USF enforcement,
  new entities with no network history, or shared owners across bidders.
\item
  \textbf{Cost outliers:} cost per location far above peers in similar
  terrain.
\end{enumerate}

\hypertarget{first-prototype}{%
\section{First prototype}\label{first-prototype}}

A BEAD subgrantee screen for two or three states that have published
awards, joining FCC map and USF history.

\hypertarget{limits}{%
\section{Limits}\label{limits}}

State award data formats vary; build-out verification needs field data.

\hypertarget{who-to-take-it-to}{%
\section{Who to take it to}\label{who-to-take-it-to}}

NTIA Office of Internet Connectivity and Growth, Commerce OIG, and state
broadband offices.

\hypertarget{sources}{%
\section{Sources}\label{sources}}

\begin{itemize}
\tightlist
\item
  \href{https://broadbandusa.ntia.gov/funding-programs/broadband-equity-access-and-deployment-bead-program}{NTIA
  BEAD program}
\item
  \href{https://broadbandmap.fcc.gov/}{FCC National Broadband Map}
\item
  \href{https://www.usaspending.gov/agency}{USAspending agency profiles
  (FY2026 outlays)}
\end{itemize}

\end{document}
